\documentclass[11pt]{letter}
\usepackage[margin=1in]{geometry}
\usepackage{hyperref}

\signature{Derek Tharp\\Department of Accounting \& Finance\\University of Southern Maine\\derek.tharp@maine.edu}

\begin{document}

\begin{letter}{Editors\\Journal of Public Economics}

\opening{Dear Professors Hendren and Kopczuk,}

I am submitting ``Quantifying the Annuity Puzzle: A Unified Lifecycle Decomposition'' for consideration at the \textit{Journal of Public Economics}.

Yaari (1965) proved that rational consumers should fully annuitize under actuarially fair pricing, yet voluntary ownership among US retirees is only 3--6\%. Six decades of research have proposed partial explanations for this gap, but the literature has lacked a unified quantitative accounting of how the main channels interact in a single model. Pashchenko (2013, \textit{JPubE}) came closest, but still predicted participation of roughly 20\%.

This paper nests nine channels within a single calibrated lifecycle framework. The results have a two-layer structure. Seven standard rational channels reduce predicted ownership from a frictionless population benchmark of 41.4\% to 18.3\%, closely matching the residual overprediction in Pashchenko (2013). Adding age-varying consumption needs---calibrated to the retirement expenditure decline documented by Aguiar and Hurst (2013)---reduces predicted ownership to 7.4\%. The full nine-channel model predicts 6.6\%, within the observed 3--6\% range.

The paper makes three contributions. First, it identifies age-varying consumption needs as a quantitatively important channel missing from prior annuity-puzzle models. Second, it provides an exact Shapley decomposition computed over all 512 channel subsets, delivering order-independent attribution and showing that pricing loads are the largest single demand suppressor. Third, it contributes directly to public-finance questions around Social Security and retirement-market design: Social Security both crowds in and crowds out private annuitization depending on the environment, group pricing (MWR = 0.90) raises predicted ownership to 43.1\%, and annuity market access has substantial welfare value for wealthy households.

I believe the paper is a good fit for the \textit{Journal of Public Economics}. It speaks directly to the journal's interests in Social Security, retirement saving, insurance-market frictions, and the design of public and employer-sponsored retirement institutions. It also directly engages with Pashchenko (2013) and related work published in \textit{JPubE} and neighboring outlets.

Replication code in Julia and all calibration data are publicly available at \url{https://github.com/DerekTharp/annuity-puzzle-model}. The manuscript has not been submitted elsewhere.

\closing{Sincerely,}

\end{letter}
\end{document}
